\documentclass[11pt]{article}
\usepackage{fontspec}
\setmainfont[
  Path=fonts/,
  Extension=.ttf,
  UprightFont=*-400,
  BoldFont=*-700
]{NotoSerifSC}
\setsansfont[
  Path=fonts/,
  Extension=.ttf,
  UprightFont=*-400,
  BoldFont=*-700
]{NotoSerifSC}
\XeTeXlinebreaklocale "zh"
\XeTeXlinebreakskip=0pt plus 1pt
\usepackage[a4paper,margin=1.70cm]{geometry}
\usepackage{amsmath,amssymb,mathtools,bm}
\usepackage{booktabs,longtable,array}
\usepackage{enumitem}
\setlength{\parindent}{2em}
\setlength{\parskip}{0.35em}
\setlist{nosep,leftmargin=2em}
\allowdisplaybreaks[3]
\numberwithin{equation}{section}
\pagestyle{plain}

\newcommand{\R}{\mathbb{R}}
\newcommand{\E}{\mathbb{E}}
\newcommand{\Pp}{\mathbb{P}}
\newcommand{\1}{\mathbf{1}}
\newcommand{\T}{\mathsf{T}}
\newcommand{\F}{\mathrm{F}}
\newcommand{\cN}{\mathcal{N}}
\newcommand{\cL}{\mathcal{L}}
\newcommand{\cS}{\mathcal{S}}
\newcommand{\cH}{\mathcal{H}}
\newcommand{\cA}{\mathcal{A}}
\newcommand{\cM}{\mathcal{M}}
\newcommand{\cT}{\mathcal{T}}
\newcommand{\diag}{\operatorname{diag}}
\newcommand{\tr}{\operatorname{tr}}
\newcommand{\Var}{\operatorname{Var}}
\newcommand{\Cov}{\operatorname{Cov}}
\newcommand{\softmax}{\operatorname{softmax}}
\newcommand{\KL}{D_{\mathrm{KL}}}
\newcommand{\sg}{\operatorname{sg}}
\begin{document}
    \begin{center}
      {\LARGE\bfseries 08\_Building Vision Models upon Heat Conduction}
    \end{center}
    \vspace{0.8em}

    \section{符号与维度}
    \begin{longtable}{>{$}l<{$} >{$}l<{$} p{10cm}}
    \toprule
    \text{符号} & \text{维度} & \text{含义}\\
    \midrule
    u(x,y,t) & \R & 连续温度或特征场。\\
        \kappa & \R_{\ge0} & 热扩散系数。\\
        \omega_x,\omega_y & \R & 二维空间频率。\\
        U^t & \R^{H\times W} & 离散视觉特征图。\\
        H_t & \R^{H\times W} & 频域热核。\\
    \bottomrule
    \end{longtable}

    所有向量默认是列向量；未特别说明时，期望同时对数据、噪声和采样时刻取值。下文只展开论文中承担方法逻辑的公式，并把所需假设写在推导发生的位置。

\section{热方程的谱解}

连续热方程
\begin{equation}
 \partial_tu(x,t)=\kappa\Delta u(x,t),\qquad \kappa>0
\end{equation}
在矩形区域上配 Neumann 边界 $\partial_nu=0$。先看一维 $x\in[0,L]$，分离变量
$u(x,t)=X(x)T(t)$：
\begin{equation}
 X(x)T'(t)=\kappa X''(x)T(t).
\end{equation}
两边除以 $\kappa XT$ 得常数 $-\lambda$，所以
\begin{equation}
 X''+\lambda X=0,\qquad T'+\kappa\lambda T=0.
\end{equation}
边界 $X'(0)=X'(L)=0$ 给出本征函数和本征值
\begin{equation}
 X_k(x)=\cos(k\pi x/L),\qquad
 \lambda_k=(k\pi/L)^2,\qquad k=0,1,\ldots
\end{equation}
时间解
\begin{equation}
 T_k(t)=e^{-\kappa\lambda_kt}T_k(0).
\end{equation}
故 Fourier--cosine 系数逐频率衰减：
\begin{equation}
 \widehat u_k(t)=e^{-\kappa(k\pi/L)^2t}\widehat u_k(0).
\end{equation}
二维本征函数是两个余弦的乘积，
\begin{equation}
 \widehat u_{k\ell}(t)=
 e^{-\kappa[(k\pi/L_x)^2+(\ell\pi/L_y)^2]t}
 \widehat u_{k\ell}(0).
\end{equation}
高频因指数中的平方频率而更快衰减，零频 $k=\ell=0$ 保持不变。

\section{离散 DCT 与 Neumann 边界}

一维网格的反射边界可写为虚拟点
$u_{-1}=u_1$、$u_N=u_{N-2}$。离散 Laplacian
\begin{equation}
 (Lu)_j=u_{j-1}-2u_j+u_{j+1}
\end{equation}
由离散余弦基向量对角化。对无限/周期直觉，代入模式
$u_j=e^{i\omega j}$：
\begin{align}
 (Lu)_j
 &=e^{i\omega j}(e^{-i\omega}-2+e^{i\omega})\\
 &=-4\sin^2(\omega/2)u_j.
\end{align}
反射边界选择实余弦模式，故 DCT 矩阵 $C$ 满足
\begin{equation}
 L=C^T\Lambda C,\qquad
 \Lambda_{kk}=-4\sin^2(\omega_k/2)\le0.
\end{equation}
离散热半群为
\begin{equation}
 e^{\kappa tL}=C^T\diag(e^{\kappa t\lambda_k})C.
\end{equation}
二维可分离情形用 $C_y\otimes C_x$ 对角化，特征值相加。

\section{显式迭代与稳定条件}

显式 Euler 离散
\begin{equation}
 u^{n+1}=u^n+\eta Lu^n=(I+\eta L)u^n.
\end{equation}
在特征模态 $k$ 上放大因子为
\begin{equation}
 g_k=1+\eta\lambda_k.
\end{equation}
稳定要求 $|g_k|\le1$ 对所有 $k$ 成立，即
\begin{equation}
 -2\le\eta\lambda_k\le0.
\end{equation}
若一维未除网格步长的 Laplacian 有 $\lambda_{\min}\ge-4$，充分条件为
$0\le\eta\le1/2$；二维五点格式最小特征值约 $-8$，充分条件变为
$\eta\le1/4$。直接使用指数谱门
$e^{\eta\lambda_k}$ 则对任意 $\eta\ge0$ 都不放大模态。

\section{能量与均值的单调性}

Neumann 边界下积分分部：
\begin{align}
 \frac d{dt}\frac12\|u(t)\|_{L^2}^2
 &=\int_\Omega u\partial_tu\,dx\\
 &=\kappa\int_\Omega u\Delta u\,dx\\
 &=-\kappa\int_\Omega\|\nabla u\|_2^2dx
 +\kappa\int_{\partial\Omega}u\partial_nu\,dS\\
 &=-\kappa\|\nabla u\|_{L^2}^2\le0.
\end{align}
总质量满足
\begin{align}
 \frac d{dt}\int_\Omega u\,dx
 &=\kappa\int_\Omega\Delta u\,dx
 =\kappa\int_{\partial\Omega}\partial_nu\,dS=0.
\end{align}
因此热扩散保留均值但降低非恒定部分能量。

离散对称半负定 $L$ 也有
\begin{equation}
 \frac d{dt}\frac12\|u\|_2^2
 =\kappa u^TLu\le0.
\end{equation}
若学习频率门 $g_k$，要保持非扩张，需约束 $|g_k|\le1$；任意可学习门可能放大高频，
此时模块不再严格对应物理热方程。

\section{频域门的梯度}

设二维 DCT 后 $\widehat X=CXC^T$，输出
\begin{equation}
 Y=C^T(G\odot\widehat X)C,
\end{equation}
其中 $C$ 正交、$G$ 是频率门。令上游梯度 $H=\partial\cL/\partial Y$，
利用 Frobenius 内积的正交不变性，
\begin{align}
 d\cL
 &=\langle H,C^T(dG\odot\widehat X+G\odot d\widehat X)C\rangle_F\\
 &=\langle CHC^T, dG\odot\widehat X+G\odot d\widehat X\rangle_F.
\end{align}
故
\begin{equation}
 \nabla_G\cL=(CHC^T)\odot\widehat X,
\end{equation}
以及
\begin{equation}
 \nabla_X\cL=C^T\left(G\odot(CHC^T)\right)C.
\end{equation}
若 $G_{k\ell}=e^{-t\lambda_{k\ell}}$ 且 $\lambda_{k\ell}\ge0$，
\begin{equation}
 \frac{\partial G_{k\ell}}{\partial t}=-\lambda_{k\ell}G_{k\ell},
\end{equation}
所以大频率对扩散时间的梯度更敏感。

\section{复杂度比较}

直接空间卷积核大小 $K\times K$、通道 $C$ 的深度卷积量为
$O(HWCK^2)$。二维快速 DCT 的量约为
\begin{equation}
 O(CHW(\log H+\log W)),
\end{equation}
再加逐频率门 $O(CHW)$。频域方法一次覆盖全局感受野；
但小分辨率或小卷积核时，变换常数和内存搬运可能使理论复杂度优势无法兑现。

\section{光滑损失的下降引理}

若 $L$ 的梯度是 $\beta$-Lipschitz：
\begin{equation}
 \|\nabla L(x)-\nabla L(y)\|
 \le\beta\|x-y\|,
\end{equation}
则沿线段积分
\begin{align}
 L(y)-L(x)
 &=\int_0^1\nabla L(x+t(y-x))^T(y-x)dt\\
 &=\nabla L(x)^T(y-x)+R,
\end{align}
余项满足
\begin{align}
 |R|
 &\le\int_0^1\beta t\|y-x\|^2dt
 =\frac\beta2\|y-x\|^2.
\end{align}
故下降引理
\begin{equation}
 L(y)\le L(x)+\nabla L(x)^T(y-x)
 +\frac\beta2\|y-x\|^2.
\end{equation}
取梯度步 $y=x-\eta\nabla L(x)$：
\begin{equation}
 L(y)\le L(x)-\eta\left(1-\frac{\beta\eta}{2}\right)
 \|\nabla L(x)\|^2.
\end{equation}
当 $0<\eta<2/\beta$ 时每个精确梯度步下降；深网中 $\beta$ 未知且随位置变化，
这仍解释了学习率过大造成不稳定的局部机制。

\section{二阶近似与曲率方向}

在 $x$ 附近 Taylor 展开
\begin{equation}
 L(x+\Delta)=L(x)+g^T\Delta
 +\frac12\Delta^TH\Delta+O(\|\Delta\|^3).
\end{equation}
若 $H=U\diag(\lambda_i)U^T$，在本征方向 $u_i$ 上步长 $a$ 的变化
\begin{equation}
 \Delta L\approx a g_i+\frac12\lambda_i a^2.
\end{equation}
$\lambda_i>0$ 为局部碗形，$\lambda_i<0$ 是负曲率逃离鞍点方向。
加阻尼的 Newton 步
\begin{equation}
 \Delta=-(H+\lambda I)^{-1}g
\end{equation}
在特征方向缩放为 $-g_i/(\lambda_i+\lambda)$，避免接近零或负特征值造成巨大步长。

\section{链式 Jacobian 与条件数}

复合映射 $y=f_L\circ\cdots\circ f_1(x)$ 的 Jacobian
\begin{equation}
 J_{y\leftarrow x}=J_LJ_{L-1}\cdots J_1.
\end{equation}
谱范数上界
\begin{equation}
 \|J_{y\leftarrow x}\|_2\le\prod_{\ell=1}^{L}\|J_\ell\|_2.
\end{equation}
若各层最小奇异值存在，
\begin{equation}
 \sigma_{\min}(J_{y\leftarrow x})
 \ge\prod_\ell\sigma_{\min}(J_\ell).
\end{equation}
因此条件数至多乘法累积：
\begin{equation}
 \kappa(J_{y\leftarrow x})
 \le\prod_\ell\kappa(J_\ell).
\end{equation}
低维映射、递归模块和物理传播都面临同一问题：表示存在不代表反向梯度数值良态。

\section{随机梯度的偏差与方差}

设样本梯度 $g(\theta;\xi)$ 满足
\begin{equation}
 \E_\xi g(\theta;\xi)=\nabla L(\theta),
 \qquad \E\|g-\nabla L\|^2\le\sigma^2.
\end{equation}
独立批量 $B$ 的平均
\begin{equation}
 \bar g_B=\frac1B\sum_{i=1}^{B}g_i
\end{equation}
满足
\begin{equation}
 \E\bar g_B=\nabla L,\qquad
 \E\|\bar g_B-\nabla L\|^2\le\frac{\sigma^2}{B}.
\end{equation}
若样本相关，协方差项出现：
\begin{equation}
 \Var(\bar g_B)=\frac1{B^2}\left[
 \sum_i\Var(g_i)+\sum_{i\ne j}\Cov(g_i,g_j)
 \right].
\end{equation}
扩大批量只有在交叉协方差不主导时才近似按 $1/B$ 降噪。

若使用有偏估计 $\E\widehat g=\nabla L+b$，均方误差分解
\begin{equation}
 \E\|\widehat g-\nabla L\|^2
 =\|b\|^2+\tr\Cov(\widehat g).
\end{equation}
直通估计、停止梯度目标和截断反传通常以可控偏差换低方差或低内存。

\section{Adam 的尺度归一化}

逐坐标 Adam 更新
\begin{align}
 m_t&=\beta_1m_{t-1}+(1-\beta_1)g_t,\\
 v_t&=\beta_2v_{t-1}+(1-\beta_2)g_t^2,\\
 \widehat m_t&=m_t/(1-\beta_1^t),\\
 \widehat v_t&=v_t/(1-\beta_2^t),\\
 \theta_{t+1}&=\theta_t-\eta
 \frac{\widehat m_t}{\sqrt{\widehat v_t}+\epsilon}.
\end{align}
若梯度长期乘常数 $c>0$，忽略 $\epsilon$ 时
$\widehat m$ 乘 $c$、$\sqrt{\widehat v}$ 也乘 $c$，更新近似不变。
但瞬态、符号变化和 $\epsilon$ 会破坏精确尺度不变性。

AdamW 把权重衰减解耦：
\begin{equation}
 \theta_{t+1}=(1-\eta\lambda)\theta_t
 -\eta\frac{\widehat m_t}{\sqrt{\widehat v_t}+\epsilon}.
\end{equation}
这不同于把 $\lambda\theta$ 加入梯度后再经过自适应分母；后者会按坐标二阶矩缩放正则化。

\section{范数约束与投影梯度}

约束集合 $\mathcal C$ 上的投影梯度
\begin{equation}
 \theta^+=\Pi_{\mathcal C}(\theta-\eta g),
 \qquad
 \Pi_{\mathcal C}(x)=\arg\min_{y\in\mathcal C}\|x-y\|^2.
\end{equation}
对欧氏球 $\mathcal C=\{x:\|x\|\le R\}$，
\begin{equation}
 \Pi_{\mathcal C}(x)=
 \begin{cases}
 x,&\|x\|\le R,\\
 Rx/\|x\|,&\|x\|>R.
 \end{cases}
\end{equation}
对流形约束，局部一阶步先投影到切空间，再用 retraction 回到流形；
仅把参数写成 $g(z)$ 相当于用参数化隐式满足约束，但会引入 $J_g^TJ_g$ 的坐标度量。

\section{Lipschitz 误差的层间传播}

真实模块 $f_\ell$、近似模块 $\widehat f_\ell$ 满足
\begin{equation}
 \|f_\ell(x)-\widehat f_\ell(x)\|\le\epsilon_\ell,
 \qquad
 \|f_\ell(x)-f_\ell(y)\|\le L_\ell\|x-y\|.
\end{equation}
令真实、近似中间状态差 $e_\ell$。三角不等式给出
\begin{align}
 e_{\ell+1}
 &=\|f_\ell(h_\ell)-\widehat f_\ell(\widehat h_\ell)\|\\
 &\le L_\ell e_\ell+\epsilon_\ell.
\end{align}
递归展开
\begin{equation}
 e_L\le
 \left(\prod_{j=0}^{L-1}L_j\right)e_0
 +\sum_{\ell=0}^{L-1}
 \left(\prod_{j=\ell+1}^{L-1}L_j\right)\epsilon_\ell.
\end{equation}
早期层误差乘更多后续 Lipschitz 因子，所以量化、低秩近似或局部训练的相同单层误差，
放在不同深度并不等价。

\section{复杂度必须区分四个量}

参数量、FLOPs、激活内存和通信量是不同指标。以线性层
$Y=XW$、$X\in\R^{B\times n}$、$W\in\R^{n\times m}$ 为例：
\begin{align}
 \text{参数量}&=nm,\\
 \text{前向乘加量}&\asymp Bnm,\\
 \text{主要激活量}&=B(n+m),\\
 \text{数据并行梯度通信}&\asymp nm.
\end{align}
低秩因子可同时减少参数与乘法，但动态稀疏若实现仍做稠密 kernel，参数/FLOPs 公式不会自动转化为墙钟加速。
冻结参数减少梯度和优化器状态，却不一定减少前向计算或输入梯度计算。

\end{document}
